% Immutable cited-input supplement. Target proofs remain embargoed.
% Target: night:P4.26; tag 02WY.
% Frozen source SHA256 74497bfe327fae25c9cdcb17334b3bf42a9793a9df7c0a4404d97568f2f760f0.
% CompendiumRef is resolved against the cited S6_09 source edition, not silently identified with the current series statement.

% Exact cited input: series:II / night:P4.12
\begin{proposition}[P4.12]\hypertarget{tag-02WJ}{}\label{night:P4.12}\label{stmt:P4.12}


In the standing torus-fibration setting of this section, let $f^{-1}(b)=mS$
have irreducible reduced support. Then
\[
 \omega_S^{\otimes m}\simeq\OO_S,\qquad
 \chi(\OO_S)=0,\qquad h^1(\OO_S)\in\{1,2\}.
\]
Every normal such support is smooth. For a singular, necessarily
nonnormal support, a minimal resolved model of its normalization
has Kodaira dimension $-\infty$.
\end{proposition}

% Exact cited input: series:II / night:P4.13
\begin{proposition}[P4.13]\hypertarget{tag-02WK}{}\label{night:P4.13}\label{stmt:P4.13}


Let a proper holomorphic degeneration of two-dimensional complex
tori have smooth total threefold and central fibre $mS$ with
irreducible reduced support. If $A$ is the monodromy on $H_1$ of
a nearby torus, then
\[
 m e(S_{\mathrm{reg}})=\det(I-A^m),\qquad
 0\le e(S_{\mathrm{reg}})\le16/m.
\]
This assertion requires no Betti-number hypothesis on a compactification.
\end{proposition}

% Exact cited input: series:II / night:P4.21
\begin{proposition}[P4.21]\hypertarget{tag-02WT}{}\label{night:P4.21}\label{stmt:P4.21}


In the standing torus-fibration setting of this section, an irreducible singular support $S$
with smooth rational normalization $T$ has $e(S)>0$.
\end{proposition}

% Exact cited input: series:II / night:P4.23
\begin{proposition}[P4.23]\hypertarget{tag-02WV}{}\label{night:P4.23}\label{stmt:P4.23}


In the standing torus-fibration setting of this section, let an irreducible singular support
$S$ have smooth normalization $T$. If its minimal model $T_0$
belongs to class VII and $b_2(T_0)>0$, then $e(S)>0$.
\end{proposition}

% Exact cited input: series:II / night:P4.25
\begin{proposition}[P4.25]\hypertarget{tag-02WX}{}\label{night:P4.25}\label{stmt:P4.25}


Every semi-log-canonical reduced fibre support in
Part~4.2 has normal crossings. In particular, let
$\rho:T\to N\to S$ be a resolution of the normalization
of an irreducible support, and let $D$ be the integral
divisor defined by
$\omega_T(D)\simeq\rho^*\omega_S$. When this actual
divisor is reduced and has normal crossings, $S$
has normal crossings, its normalization is smooth,
and P4.6--P4.8 apply.
\end{proposition}

% Exact cited input: series:II / night:P4.6
\begin{proposition}[P4.6]\hypertarget{tag-02XY}{}\label{night:P4.6}\label{stmt:P4.6}


Let $X$ be a compact complex threefold with $b_2(X)=0$, and let $F$ be a
reduced compact normal-crossings divisor. We allow different local
branches of one global irreducible component to meet. If $T(F)$ is the
number of ordinary triple points, then
\[
 e(F)=T(F),\qquad T(F)\text{ is even}.
\]
This assertion concerns the reduced support, regardless of the
multiplicities with which it occurs in a fibre.
\end{proposition}

% Exact cited input: series:II / night:P4.7
\begin{proposition}[P4.7]\hypertarget{tag-02XZ}{}\label{night:P4.7}\label{stmt:P4.7}


Let $f:X\to B$ be a proper surjective holomorphic map with connected
fibres from a compact complex threefold satisfying $b_1(X)=b_2(X)=0$
to a smooth compact curve. Suppose all reduced fibre supports have
normal crossings. Then
\[
 B=\Pp^1,\qquad
 c_3[X]=\sum_{b\in B}T((f^{-1}(b))_{\mathrm{red}})\in\{0,2\},\qquad
 b_3(X)=2-c_3[X].
\]
The numerical conclusion does not require the general fibre to be a
torus.
\end{proposition}

% Exact cited input: series:II / night:P4.8
\begin{proposition}[P4.8]\hypertarget{tag-02Y0}{}\label{night:P4.8}\label{stmt:P4.8}


For a torus fibration in P4.7 whose reduced supports have normal
crossings, an irreducible singular support with
rational normalization has Euler number two; every other irreducible
support has Euler number zero. In the zero case the supports have the
forms stated in the prompt: a smooth support has a finite unramified
torus cover; a singular irreducible support has smooth minimal
elliptic-ruled or Hopf normalization; a reducible support has smooth
minimal Hopf or non-Kähler elliptic normalization components, with
smooth elliptic double curves and finite unramified branch maps.
Thus $c_3[X]=0$ precisely when every fibre belongs to these zero cases.
\end{proposition}

% Exact cited input: series:II / night:P4R.1
\begin{proposition}[P4R.1]\hypertarget{tag-02Y2}{}\label{night:P4R.1}\label{stmt:P4R.1}


For a proper holomorphic degeneration of complex two-tori with
smooth total threefold and central fibre $mS$, where $S$ is
irreducible and reduced, the group $\pi_1(S)$ has an abelian
normal subgroup with cyclic quotient of order dividing $m$.
\end{proposition}

% Exact cited input: series:II / night:P4R.6
\begin{proposition}[P4R.6]\hypertarget{tag-02YI}{}\label{night:P4R.6}\label{stmt:P4R.6}


Let $S$ be an irreducible singular support in
Part~4.2, let $\nu:N\to S$ be its normalization,
and let $\rho:T\to N\to S$ use the minimal
resolution of $N$. There is an effective integral
divisor $D$ on $T$ such that
\[
 \omega_T(D)\simeq\rho^*\omega_S,
 \qquad \operatorname{Supp}D
       =\rho^{-1}(\operatorname{Sing}S)_{\rm red}.
\]
In particular $K_T+D$ is torsion, by P4.12.
The strict-transform conductor has its usual
positive integral coefficients. Every connected
exceptional configuration over a singular point
of $N$ has strictly positive coefficients in $D$.
Every such configuration meets the strict-transform
conductor.
\end{proposition}

% Exact cited input: series:II / night:P4R.7
\begin{proposition}[P4R.7]\hypertarget{tag-02YJ}{}\label{night:P4R.7}\label{stmt:P4R.7}


In Part~4.2, every irreducible singular
support whose normalization has a
resolution birationally ruled over a
curve of genus at least two has
nonnegative Euler number. The
normalization may be singular.
\end{proposition}

% Exact cited input: series:II / night:P4R.8
\begin{proposition}[P4R.8]\hypertarget{tag-02YK}{}\label{night:P4R.8}\label{stmt:P4R.8}


In Part~4.2, every irreducible singular support
whose normalization has a resolution with Hopf
minimal model has nonnegative Euler number.
The normalization may be singular. For the
minimal normalization resolution, let $k$ be
the number of blowups from its Hopf minimal
model and $v$ the number of rational components
of the pole divisor. Then
\[
 e(S_{\rm reg})=k-v\ge0,\qquad
 e(\operatorname{Sing}S)\ge0.
\]
\end{proposition}

% Exact cited input: series:II / night:P4R.9
\begin{proposition}[P4R.9]\hypertarget{tag-02YL}{}\label{night:P4R.9}\label{stmt:P4R.9}


In Part~4.2, every irreducible singular
support whose normalization has a
resolution birationally ruled over an
elliptic curve has nonnegative Euler
number. When the horizontal part of
the pole divisor on the minimal ruled
model is reduced, the normalization
is smooth and minimal and the support
has Euler number zero.
\end{proposition}

